% !TEX root = main.tex
% 由 GearPro_项目进展情况小结.md 转换；严格模仿 MineTemplate v2 格式与工科实训规范。

\begin{center}
{\zihao{3}\bfseries GearPro 齿轮外观缺陷检测系统：项目进展情况小结\par}
\end{center}

\section*{项目阶段进展简要概括}

针对齿轮表面微细划痕与局部齿部缺损的质检难题，本项目开发了一套流水线在线辅助人工复核原型系统。在指导教师指导下，团队自主采购定制 50 个真实齿轮，采集多视角图像并建立了防数据泄漏的“近重复组”划分机制与涵盖正常、划痕、三种细分缺齿及困难样本的 6 种标签体系。针对两类缺陷的形态差异，算法历经 6 大轮实验演进，确立了齿轮 ROI 提取加局部缺陷识别的 2 级级联流水线、全局 CNN 分类与局部 P2-YOLO 目标检测的异构融合架构；通过验证集最小化负对数似然进行温度缩放校准，利用 TPE 贝叶斯优化确定集成权重，依据经验奈曼–皮尔逊准则确定工作阈值，并在决策层采用“或规则”合并报警。在 150 张离线测试图上，系统取得 88.73\% 召回率（63/71）与 22.78\% 误报率（18/79），较早期统一基线净增 9 张检出、减少 15 张误报。硬件上完成了第一代整机实物搭建，实现了平皮带逐件送料、透明亚克力载台无遮挡采光、步进电机平滑加减速旋转分拣、上下双工业相机与漫射照明，集成了 STM32 串口协议与 Jetson Orin NX 本地推理（定位模型 TensorRT FP16 耗时压缩至 26.7 ms），打通了全流程闭环。目前系统仍存在未搭建基于光度立体视觉的多光源成像系统、数据未进行物理工件 ID 严格隔离、数据集规模与划分待优化及边缘部署端探索较少等局限，后续将深入优化。

\section*{一、 项目背景与实际质检问题}

在机械传动系统中，齿轮是基础的核心传动构件。在机械加工、热处理与流转装配中，齿面易受工具划擦或外力磕碰，产生表面微细划痕与局部齿部缺损（含齿侧碰伤与齿角崩缺）。若带缺陷齿轮流入装配，易诱发啮合噪声加剧与应力集中。

车间调研表明，传统质检主要依靠工人肉眼目检，存在劳动强度大、连续作业易疲劳漏检以及检验员尺度不一等局限。因此，我们希望利用机器视觉技术开发在线筛查原型系统，在流水线上自动筛查疑似缺陷齿轮供人工重点复核，降低漏检率并减轻工人负担。

需要说明的是，本项目现阶段定位为\textbf{辅助人工复核的筛查工具}。由于目前尚未建立齿轮疲劳寿命与机械强度的定量验收标准，系统判定的核心依据是图像上肉眼可见的异常划痕与形貌破坏，识别出“疑似缺陷”并不等同于对齿轮机械安全性能下终检定论。

在团队分工上，由 3 名大三工科学生协作开展：本人负责缺陷数据集整理标注、算法设计训练、多模型异构融合及上位机检测软件与 Web 监控界面的编写；机械与光学成员负责第一代台架机械结构设计、透明载台加工及相机光源选型打光；电气控制成员负责基于 STM32 的电控编程、步进电机驱动与平滑加减速控制及串口协议联调。目前第一代实物台架已搭建完成并打通全流程闭环。

\section*{二、 文献研读与技术路线启发}

针对齿轮质检中遇到的金属反光强、微细划痕难提取及立体齿廓多角度遮挡难题，我们在指导教师指导下研读了国内外相关文献，从中获得了重要技术启发。

\subsection*{1. 高反光金属打光与成像方案的启发}
通过对文献的阅读，我们了解到高反光金属工件的光照设计直接决定了缺陷在图像中的可见度。Saiz 等人在研究反光金属零件外观检测时指出，金属镜面反射亮斑与暗影极易掩盖微小缺陷；Wu 等人在复杂金属构件检测中也指出，受控的多角度漫射照明对于压制局部高光反射至关重要。齿轮具有曲面渐开线齿廓，直射光在齿顶与齿根极易形成强烈反射亮斑与死角阴影，遮蔽微细划痕。结合文献建议与试验对比，我们在台架中选用了高亮度白色漫射光源配合导光罩与上下双相机布局，有效弱化了金属强反光对浅划痕的遮蔽，改善了成像质量。

\subsection*{2. 微细划痕特征保留的方法借鉴}
在算法特征提取方面，通过对多尺度网络文献的学习，我们发现常规卷积神经网络在多次下采样后，微小目标的高分辨率空间特征容易丢失。Lin 等人提出的特征金字塔网络（FPN）通过横向连接融合深层语义与浅层高分辨率特征；Liu 等人在带钢缺陷分类中提出的 ConCNN 结构也证实保留高分辨率分支有助于微小特征捕获。齿轮微细划痕宽度通常仅数个像素，若直接多层下采样极易丢失特征。借鉴文献思想，我们在目标检测模型中引入了专门保留浅层高分辨率特征的 \textbf{P2 检测层}，专门负责捕获细小划痕，为提高微小缺陷检出率提供了结构基础。

\subsection*{3. 旋转多角度拍摄与连续图像判定的文献启发}
我们重点研读了 Allam 等人关于汽车齿轮表面损伤视觉检测的研究。作者采用多视角相机配合夹持旋转机构逐齿扫描齿轮，指出齿面损伤往往只在特定旋转角度下可见；更为关键的是，作者对比了单图判定与连续多帧判定规则，表明单帧判断极易受瞬间反光噪声干扰产生误报，而采用“连续多帧均检出才报警”的滑动窗口规则，能在保证缺陷检出的同时显著降低正常齿轮的误报率。这启发我们：齿轮必须在旋转运动中多角度拍摄才能完整覆盖全部齿面；同时算法后续必须结合旋转连续帧的时序信息进行综合判定，才能有效过滤偶然反光带来的单帧误报。

\section*{三、 数据构建、算法演进与决策依据}

软件与算法是本项目前期投入精力最多的核心模块。在缺乏现成齿轮缺陷公开数据集的情况下，我们结合自建数据，自主设计了防泄漏划分与标签体系，并经历了 6 大轮实验演进与参数优化。

\subsection*{1. 自主设计数据集、防泄漏划分与 6 种标签体系}
我们自行采购并定制了 50 个真实齿轮（10 个正常、15 个细划痕、15 个侧部磕碰缺损、10 个齿角崩缺），利用台架上的工业相机多角度拍摄建库。

在标注规范上，我们建立并演进了逐渐严谨的 \textbf{6 种标签体系}：\texttt{normal}（正常外观）、\texttt{scratch}（微细划痕）、三种细分缺齿位置标签 \texttt{bottom}（齿根碰伤）、\texttt{oblique}（齿面斜向崩缺）、\texttt{side}（齿侧剥落），以及 \texttt{difficult}（对角度极限或弱光导致人眼极其模糊、难以准确标定边界的目标标记为困难样本，在训练集中予以剔除，避免模糊标注干扰网络梯度收敛）。

在数据集划分上，为了防止数据泄漏，我们设计了\textbf{基于“近重复组”的验证集划分方法}：由于同一齿轮旋转拍摄会产生多张角度相近的图片，若随机打乱极易导致同一工件的高度相似图像同时进入训练集与测试集，造成性能虚高。因此，我们利用 64 位感知哈希（dHash，汉明距离 $\le 2$ 合并）、文件序号邻近（$\le 3$）及长宽比推断视角（\texttt{face}, \texttt{oblique}, \texttt{side}）聚类构建“近重复组”；随后采用分层抽样平衡类别与视角分布，以“组”为单位整体划分训练集、验证集与测试集（大组保留在训练集），并在划分后交叉核验跨集 dHash 距离，严格杜绝了相似样本跨集渗透。

\subsection*{2. ROI 提取与 2 级级联检测流水线}
吸取早期背景干扰严重的教训，我们确立了“ROI 提取 + 局部缺陷识别”的 \textbf{2 级级联检测流水线}：第 1 级由轻量级定位网络快速定位整颗齿轮外接轮廓，并在宽高方向外扩 4\% 裁剪出感兴趣区域（ROI）；第 2 级将裁剪后的标准化 ROI 送入下游专项缺陷模型深入识别。外扩 4\% 既保留了边缘齿廓完整性，又排除了整帧背景杂光干扰，为后续模型稳定判别提供了规范输入。

\subsection*{3. 六大轮算法实验演进历程}
算法方案经历了从早期失败摸索到逐步成熟的 6 大轮演进历程：
\begin{enumerate}
    \item \textbf{第 1 轮（传统图像处理失败探索）}：尝试使用 Canny 边缘检测与自适应阈值分割等传统方法，受金属反光亮斑与齿廓伪边缘干扰产生大量虚报，算法完全无法分离真实缺陷与正常背景；
    \item \textbf{第 2 轮（整图端到端检测探索）}：尝试在未截取 ROI 的整图上直接训练通用 YOLO 模型，原图背景占比大且微细划痕特征微弱，模型将大量细划痕误判为正常车削刀纹，划痕漏检率极高；
    \item \textbf{第 3 轮（引入 ROI 二级流水线，探索分类与定位指标差异）}：引入齿轮 ROI 裁剪后，分别测试了图像分类网络与目标检测定位模型，发现整图分类擅长全局粗判，检测模型能给出具体坐标框，明确了两者评价维度不同，需协同设计；
    \item \textbf{第 4 轮（统一缺陷大模型受挫与反思）}：尝试将划痕与缺损合并训练“统一二分类大模型”并做多模型融合，测试发现大模型在缺损上表现尚可但微细划痕漏检依然严重，单纯增加模型数量未能解决漏检，反而倍增计算开销；
    \item \textbf{第 5 轮（分缺陷分别建模）}：针对划痕（线状弱对比）与齿部缺损（块状几何破坏）物理特性的巨大差异，确立分别建模路线——划痕子系统采用含 P2 高分辨率层的 YOLO 与分类模型融合，缺损子系统采用单类 YOLO 与分类模型融合；
    \item \textbf{第 6 轮（最终确定决策级融合与或规则）}：确立最终决策级融合（Decision-level Fusion）方案，两类缺陷子系统各自独立设阈值二值判定，最终采用“逻辑或规则（OR Rule）”合并报警——任一子系统判定为疑似缺陷即触发报警，实现分项专精、联合兜底。
\end{enumerate}

\subsection*{4. 异构模型融合、概率校准与决策参数优化机制}
针对多模型融合与判决决策，我们设计了具有充分理论依据的技术闭环：
\begin{itemize}
    \item \textbf{多模型异构融合架构}：单类缺陷内采用“全局 CNN 分类（EfficientNet-B0、ResNet18 提取全图语义抗噪） + 局部目标检测（P2-YOLO 提取局部细节坐标置信度）”的异构组合，实现全局与局部的优势互补；
    \item \textbf{温度缩放（Temperature Scaling）概率校准}：异构模型输出的 logit/置信度尺度不一，直接加权会导致分数偏移。我们在融合前利用验证集样本，通过网格搜索最小化负对数似然（NLL），为每个模型独立学习温度参数 $T$（划痕三成员 $T$ 分别为 0.700、2.050、2.125；缺损三成员 $T$ 分别为 1.000、2.375、1.850），通过公式 $p_T = \sigma(\text{logit}(p)/T)$ 将各异构模型输出精准映射对齐到可比概率尺度；
    \item \textbf{TPE 贝叶斯优化确定集成权重（核心重点）}：在得到校准概率后，为了避免人工经验设定加权系数，我们引入基于树状 Parzen 估计器（TPE）的贝叶斯优化算法，在验证集上全局寻优确定最优加权权重。最终确定划痕子系统中目标检测模型权重为 0.75、两分类模型各为 0.125（充分发挥 P2 检测层对微小划痕的定位捕捉能力）；齿部缺损子系统中目标检测权重为 0.50、两分类模型各为 0.25；
    \item \textbf{经验奈曼–皮尔逊准则（Empirical NP Rule）确定判定阈值}：在工业复核场景中必须严格约束正常件误报成本。我们在验证集上基于经验奈曼–皮尔逊准则，在设定假阳性率上限（$\text{FPR} \le 0.20$）的前提下最大化召回率，并兼顾各位置组最低召回率，最终选定划痕判定阈值 $\theta \approx 0.3003$、缺损判定阈值 $\theta \approx 0.3413$；
    \item \textbf{离线测试集实测结果}：在严格锁定的 150 张离线测试集（31 划痕、44 缺损、4 兼具，共 71 真实缺陷、79 正常）上，最终第 6 轮方案成功检出 63/71 张缺陷图，召回率达到 88.73\%；误报 18/79 张正常图，误报率为 22.78\%。相比第 4 轮早期统一分类基线（检出 54/71、误报 33/79），\textbf{正确检出数净增 9 张，误报数减少 15 张，误报率自 41.77\% 下降至 22.78\%}，充分验证了分缺陷建模、温度校准与贝叶斯权重融合的有效性。
\end{itemize}

\section*{四、 第一代整机硬件设计与闭环联调进展}

为了检验算法在真实设备上的执行表现，团队完成了第一代整机实物台架的搭建与调试，打通了从机械传送、光学成像、边缘计算到电气分拣的全部物理与通信链路。

\subsection*{1. 机械结构、逐件传送与透明检测工位}
团队完成了整机机械构件的加工与装配：
\begin{itemize}
    \item \textbf{逐件输送机构}：设计装配了步进电机驱动的平皮带输送机构，满足齿轮“一次只送入一个工件”逐件连续输送至光学检测区的工作要求；
    \item \textbf{透明检测工位载台}：在视觉检测工位专门设计加工了高透光透明亚克力托盘载台，在稳定承托齿轮的同时为上下两台工业相机提供无物理遮挡的采光通道，便于两路相机同时观察齿轮正反两面。
\end{itemize}

\subsection*{2. 旋转分拣平台与执行机构}
在分拣工位设计搭建了由步进电机驱动的旋转式分拣平台，负责在接收到上位机判决指令后将齿轮送入对应的分类收集工位（合格品工位 / 疑似缺陷工位）。针对机械转动过程中的惯性冲击难题，在运动控制上编写并实现了\textbf{步进电机梯形平滑加减速算法}，精确控制启动、匀速与制动时序，有效避免了因启停加速度过大将齿轮甩离载台，保障了机械运转的平稳性。

\subsection*{3. 光学与机器视觉采集硬件}
在机械支架上完成了上下双工业相机的空间布局与实物安装：上方相机垂直对准齿轮上端面与齿顶轮廓，侧下方相机对准齿面渐开线与齿根，选用原图分辨率为 2560 $\times$ 1440 的工业面阵相机；光源配置了高亮度白色漫射光源，有效抑制了金属表面的强光反射亮斑。

\subsection*{4. 电气控制与微控制器（下位机 STM32）}
选用 STM32 单片机作为底层电控核心：
\begin{itemize}
    \item 编写固件实现了对输送带启停运行速度的控制，以及对旋转分拣平台步进电机的精准位置环加减速控制；
    \item STM32 与上位机通过\textbf{硬件串口（UART，115200 波特率）}实现物理连通与协议解析，可稳定接收并执行两类单字节指令：\texttt{01} 为未检出缺陷（PASS），分拣盘旋转至合格工位放行；\texttt{02} 为疑似缺陷（REJECT），分拣盘旋转至缺陷工位分选。
\end{itemize}

\subsection*{5. 边缘计算部署与全流程闭环打通}
选用并部署了 \textbf{NVIDIA Jetson Orin NX} 嵌入式计算模块，所有检测算法均在本地端独立运行。实验中，我们将齿轮定位模型成功编译转换为 \textbf{TensorRT FP16 引擎}，单环节推理耗时（P95）从原始 PyTorch 的 65.0 ms 压缩降至 26.7 ms，显著降低了端到端时延；上位机配套开发了 Web 监控控制台，现场实时显示抓拍画面、缺陷红框坐标与运行状态。

经实地联调，系统成功打通并连通了：\textbf{输送带送料 $\rightarrow$ 工业相机采集 $\rightarrow$ Orin NX 定位与缺陷判定 $\rightarrow$ Web 监控显示 $\rightarrow$ 串口通信 $\rightarrow$ STM32 控制步进电机旋转分拣}的全套物理执行与控制闭环，验证了整机方案的工程可行性。

\section*{五、 当前存在的不足与未来改进方向}

结合指导老师审读意见与实际台架联调体会，我们对当前系统的局限性进行了客观梳理。

\subsection*{1. 当前存在的不足与工程局限}
\begin{enumerate}
    \item \textbf{未来得及搭建基于光度立体视觉（Photometric Stereo）的多光源成像系统}：通过研读文献（如 Saiz 等与 Wu 等关于金属零件检测的研究），我们了解到抑制高反光最有效的物理手段是搭建光度立体视觉系统：在工件静止下通过 4 方向可控 LED 快速交替点亮，拍摄多幅不同光照图像并解耦计算出原始纹理图（Texture）、梯度幅值范围图（Range）以及表面曲率图（Curvature），合成伪彩色多通道输入送入神经网络，从而在物理成像源头将真实几何凹凸与单纯反射亮斑彻底分离。本项目第一代实物台架受限于实训周期与硬件条件，仅采用了稳态白色漫射光源，尚未来得及研制基于单片机频闪同步触发的四方向光度立体硬件与采集系统；
    \item \textbf{测试数据尚未进行物理工件 ID 严格隔离}：虽然我们采用了基于 dHash 近重复聚类的划分算法，但自建数据集目前仍是以图像编号进行管理的。50 个真实齿轮的多角度照片尚未在底层数据库中建立严格的实体工件 ID 追溯机制，现有的 88.73\% 召回率是在这批测试图片上的表现，面对未见过的全新批次齿轮时的泛化能力仍需补充独立测试证据；
    \item \textbf{数据集规模有限且标注划分有待优化}：目前仅有 50 个物理齿轮样件，缺陷形态多样性与视角连续性仍显不足；后续需采集更多批次齿轮扩充数据规模，并建立更严格的基于实体 ID 隔离的划分标准与更精细的形貌标注规范；
    \item \textbf{边缘部署端探索较为初步}：目前仅对齿轮定位模型完成了 TensorRT FP16 的基础加速验证，核心缺陷识别的多模型融合系统仍基于 PyTorch 运行，整体部署端的技术探索较少；尚未对多模型流水线实施轻量化剪枝、模型蒸馏、INT8 精度量化以及硬件长时运行温升与功耗做深入系统调优。
\end{enumerate}

\subsection*{2. 下一阶段重点推进计划}
针对上述薄弱环节，团队制定了下一阶段的推进计划：
\begin{enumerate}
    \item \textbf{探索搭建多光源光度立体硬件模组}：参考 Saiz 等人方案，研制基于 STM32 微秒级频闪同步控制的四方向可控 LED 光源模组，算法端开发纹理、梯度与曲率图的解耦融合接口，从成像源头抑制反光干扰；
    \item \textbf{扩充数据集并开展独立工件盲测}：采购更多批次全新齿轮，建立严格的物理实体 ID 数据库，实施基于工件实体完全隔离的训练、验证与盲测基准，客观评估新工件泛化能力；
    \item \textbf{深化 Jetson Orin NX 边缘端全面优化部署}：探索多模型联合剪枝、知识蒸馏与 INT8 权重量化，优化端到端异步多线程调度流水线，实测长时运转温升稳定性与连续生产节拍。
\end{enumerate}

\section*{实训小结与展望}

通过这段时间的科研与工程实训，我们在指导教师引导下查阅文献学习了机器视觉打光反光抑制、微小缺陷特征保留与旋转多角度判定的基本思路。团队完成了 50 个真实齿轮的多角度建库与 6 种标签体系梳理，历经 6 轮算法演进，确定了 ROI 二级级联流水线、多模型异构融合、温度缩放概率校准、TPE 贝叶斯权重寻优、经验 NP 准则阈值与决策级或规则合并的技术路线，在离线测试集上取得了 88.73\% 的召回率。同时，团队完成了第一代整机实物台架搭建，打通了输送带送料、透明载台无遮挡采光、步进电机平滑分拣、STM32 串口协议与 Jetson Orin NX 本地推理的完整物理闭环。尽管系统在光度立体成像系统搭建、工件 ID 严格隔离、数据集优化与边缘深度部署上仍有待进一步完善，但阶段性的软硬件验证为后续工作打下了扎实的基础。后续我们将针对薄弱环节深入测试，争取使系统更加稳定实用。
